\section{Primitive-recursive analytic infrastructure}\label{sec:representation}

This section records enough numerical detail to make the later proofs genuine arithmetic schemas.  The constants are not black boxes: each is generated by a bounded loop with a rational enclosure.

\subsection{Integers, rationals, dyadics, and intervals}

An integer is coded by a sign bit and a natural magnitude.  A rational is a pair $(a,b)$ with $b>0$; equality and order are cross-multiplied integer relations.  Reduction to lowest terms is optional.  Addition, multiplication, absolute value, integer powers, finite sums and products, and dyadic rounding are primitive recursive.  A rational interval is a pair $[l,u]$ with $l\le u$.  The four-endpoint rules for interval multiplication and the reversed-endpoint rule for positive reciprocals are finite case distinctions.

Given a positive rational $q$ and precision $p$, integer bisection finds $m$ such that
\[
  \frac{m^2}{2^{2p}}\le q<\frac{(m+1)^2}{2^{2p}}.
\]
Hence $[m2^{-p},(m+1)2^{-p}]$ encloses $\sqrt q$.  The loop is bounded by an explicit power of two, and its correctness is quantifier-free induction on the search interval.  The same procedure gives an inverse-square-root name once a positive lower bound is supplied.

\subsection{A certified name for \texorpdfstring{$\pi$}{pi}}

For $0\le z\le1$, the finite identity
\[
  \frac{1}{1+t^2}=\sum_{j=0}^{N-1}(-1)^j t^{2j}
      +(-1)^N\frac{t^{2N}}{1+t^2}
\]
and integration from $0$ to $z$ give
\begin{equation}\label{eq:atan-bound}
 \left|\arctan z-\sum_{j=0}^{N-1}(-1)^j\frac{z^{2j+1}}{2j+1}\right|
 \le \frac{z^{2N+1}}{2N+1}.
\end{equation}
Machin's identity
\begin{equation}\label{eq:machin}
 \pi=16\arctan(1/5)-4\arctan(1/239)
\end{equation}
is an exact tangent-addition calculation over the rationals, with the quadrant fixed by the elementary bounds $3<\pi<4$.  Equations \eqref{eq:atan-bound}--\eqref{eq:machin} therefore produce a primitive-recursive interval $\pi_p$ of width at most $2^{-p}$.  The executable certificate uses 170 and 55 terms, then rounds outward to a dyadic interval.  This is far more precision than the displayed numerical table requires.

\subsection{Exponential and logarithm}

For $0\le y\le1$, the alternating sum
\[
 E_N^-(y)=\sum_{j=0}^{N-1}\frac{(-y)^j}{j!}
\]
brackets $e^{-y}$ with its next term:
\begin{equation}\label{eq:exp-alt}
  \min\left(E_N^-(y),E_N^-(y)+\frac{(-y)^N}{N!}\right)
  \le e^{-y}\le
  \max\left(E_N^-(y),E_N^-(y)+\frac{(-y)^N}{N!}\right).
\end{equation}
For $r\ge0$, choose a power of two $m\ge r$, compute an enclosure of $e^{-r/m}$ by \eqref{eq:exp-alt}, and raise the interval to the $m$th power by repeated squaring.  Positive exponentials are reciprocals or use the ordinary positive Taylor sum with the Lagrange remainder.  The computation, its width bound, and monotonicity are all primitive recursive.

For $x>0$, first choose an integer $k$ with $1\le 2^{-k}x<2$, and write
\[
 \log x=k\log2+2\sum_{j=0}^{N-1}
   \frac{u^{2j+1}}{2j+1}+R_N,
 \qquad u=\frac{2^{-k}x-1}{2^{-k}x+1}.
\]
Since $|u|\le1/3$ on $[1,2]$,
\begin{equation}\label{eq:log-tail}
 |R_N|\le \frac{2|u|^{2N+1}}{(2N+1)(1-u^2)}.
\end{equation}
The same series at $x=2$ gives $\log2$.  Thus logarithms on intervals separated from zero have uniform primitive-recursive names.  All later uses of $\log$ state the required positive lower bound.

\subsection{Trigonometric functions}

Range reduction modulo a rational enclosure of $2\pi$, followed by the alternating sine and cosine sums,
\[
 \sin x=\sum_{j<N}\frac{(-1)^j x^{2j+1}}{(2j+1)!}+R_N^{\sin},
 \qquad
 \cos x=\sum_{j<N}\frac{(-1)^j x^{2j}}{(2j)!}+R_N^{\cos},
\]
with the next-term bounds on $|x|\le\pi/4$, gives primitive-recursive trigonometric names.  Exact identities such as $\sin^2x+\cos^2x=1$ are proved first for finite Taylor jets with a controlled remainder or, more economically, by the finite ODE uniqueness theorem from the calculus packet.  Only the Box--Muller section needs these functions.

\subsection{Finite integration contracts}

Let $g$ be represented on $[a,b]$ with an evaluator error $2^{-r}$ and a modulus of uniform continuity $\omega$.  A midpoint mesh with $M$ cells and evaluation precision $r$ produces a rational sum $S_{M,r}$.  Choosing $M$ so that $\omega(M)\le2^{-p-2}/(b-a)$ and $r\ge p+2+\lceil\log_2 M\rceil$ gives
\[
  \left|S_{M,r}-\int_a^b g(x)\dd x\right|\le2^{-p}.
\]
When derivative bounds are available, the midpoint, trapezoidal, or Simpson remainder yields a smaller primitive-recursive mesh.  Improper integrals add a tail producer $T(p)$ satisfying
\[
  \int_{|x|>T(p)} |g(x)|\dd x\le2^{-p}.
\]
All substitutions used below are performed first on finite tagged sums.  The change in mesh and Jacobian is a rational identity; the limit is discharged by the displayed moduli.  This avoids any appeal to a general measure-theoretic change-of-variables theorem inside arithmetic.

\subsection{The gamma input and Gaussian normalization}

The gamma packet provides, for every positive rational $a$ with an apartness witness, a fast name for
\[
 \Gamma(a)=\int_0^\infty t^{a-1}e^{-t}\dd t,
\]
including the functional equation and the certified identity
\begin{equation}\label{eq:gamma-half}
 \Gamma(1/2)=\sqrt\pi.
\end{equation}
For completeness, the Gaussian application is only the positive substitution $t=x^2/2$:
\begin{align*}
 2\int_0^\infty e^{-x^2/2}\dd x
 &=\sqrt2\int_0^\infty t^{-1/2}e^{-t}\dd t\\
 &=\sqrt2\,\Gamma(1/2)=\sqrt{2\pi}.
\end{align*}
Every equality here is an equality of fast rational names.  At a requested precision, truncate both integrals using exponential tails, carry out the substitution on a rational mesh separated from zero, and absorb the small initial interval $[0,\varepsilon]$ by its length bound.  Thus we obtain the normalization contract
\begin{equation}\label{eq:gaussian-normalization}
 \int_{-\infty}^{\infty}\frac{1}{\sqrt{2\pi}}e^{-x^2/2}\dd x=1.
\end{equation}

\subsection{Gaussian tail cutoffs}

The basic tail bound needed to justify every integration by parts is deliberately crude and elementary.  For $T\ge1$ and integer $m\ge0$,
\begin{equation}\label{eq:poly-gaussian-tail}
 \int_T^\infty x^m e^{-x^2/2}\dd x
 \le 2^{m+1}m!\,e^{-T/4}.
\end{equation}
Indeed $x^2/2\ge x/4$ for $x\ge1$, while $x^m\le 2^m m!e^{x/8}$, so the integrand is bounded by $2^m m!e^{-x/8}$ after harmlessly weakening constants.  Integrating the geometric/exponential majorant gives a primitive-recursive cutoff.  Later Mills inequalities provide much sharper bounds, but \eqref{eq:poly-gaussian-tail} already legitimizes polynomially weighted boundary terms.

\begin{contract}[Analytic basis used later]\label{con:analytic-basis}
For every finite precision request and every bounded rational parameter box with the displayed apartness witnesses, there are primitive-recursive procedures returning rational intervals for $\pi$, square roots, exponentials, logarithms, trigonometric functions, Gaussian densities, bounded integrals, and Gaussian improper integrals.  Their correctness predicates are bounded equations and inequalities.  Their composition laws are proved by quantifier-free induction on the finite loops and precision budgets.
\end{contract}
